\subsection{A five-element center on four states}

For a left-$N$, right-$M$ biset $S$, retain the notation
$Y=\operatorname{hocofib}(B(N,S,M)\to BN)$.  Thus $Y$ is
the classifying space of the source star with center $M$,
leaf $N$ on $S$, and one additional trivial leaf on a
singleton stalk.

\begin{lemma}[Proved: a coned P5 center on four states]
\label{CP2:A-coned-four-center-P5}
If $M=P_5$ and $1\le |S|\le4$, then $Y$ does not have
the homotopy type of $\CP^2$, for any finite $N$.
\end{lemma}
\begin{proof}
Write $K=\{e_{ij}\}$ for the minimal band and $U=SK$.
Its orbits have either one or two points.  A singleton
core orbit is a globally right-fixed state, by
Lemma~\ref{CP2:A-act-minimal-core-orbits}; write $F$ for
the set of these states.  The set $F$ is $N$ invariant.
If $|F|=1$, its point is also $N$ fixed, giving a section
$BN\to B(N,S,M)$.  Lemma~\ref{CP2:A-coned-bridge-comparisons}(1)
then kills all second squares.

Suppose $F=\varnothing$.  There is either one two-point
core orbit and at most two transient states, or two such
orbits and no transient.  In either case the two right
row actors agree on every state.  In the first case both
rows send a state to the uniquely labelled column point
of the sole orbit; in the second case this holds on each
core orbit.  Swapping the two rows of $P_5$ and fixing
$S,N$ therefore defines an involution.

Let $B$ be the set of core orbits, so $|B|\le2$, and
write $E=B(N,S,M)$, $E_U=B(N,U,M)$ and
$Y_U=\operatorname{hocofib}(E_U\to BN)$.  The regular
two-state core action category is contractible.  The
natural $N$-equivariant core comparison identifies $E_U$
with $B(N,B)$, over $BN$.  Thus $Y_U$ is the ordinary
left $N$ act cone on $B$, and all its rational second
products vanish by the binary act theorem.

Put $L=\mathbb Q[S]/\mathbb Q[U]$.  Tensor the explicit
left projective resolution of $P_5$ in
Lemma~\ref{CP2:A-order-five-corner} with $L$.  Every band
element annihilates $L$, so only its degree-one term
survives.  The semilinear lift in
Lemma~\ref{CP2:A-mixed-five-reductions} is $-1$ on that
term.  Hence the relative right-action bar is $L[1]$
with sign $-1$, as an $N$ complex.  Taking derived $N$
coinvariants preserves this comparison.  The actual
mapping-cylinder quotient $Y/Y_U\simeq\Sigma(E/E_U)$
shows that the involution is $-1$ on every rational
relative homology group.  Exactness of rational
eigenspaces makes $H_4(Y_U)\to H_4(Y)^+$ surjective.
Under a projective-plane hypothesis the second class
spans an eigenspace and its square has sign $+1$.
The square pairs to zero with $H_4(Y)^-$ by its sign,
and with $H_4(Y)^+$ because all products on $Y_U$ vanish.
This contradicts the projective-plane square.

It remains that $|F|\ge2$.  If every core orbit is a
singleton, put $X=\mathbb Q[S]$ and $V=\mathbb Q[F]$.
Exactly the $N$-equivariant calculation in
Proposition~\ref{CP2:A-coned-three-center-P5}, whose
calculation does not bound the number of transient states,
gives the projective tensor complex
\[
 V\oplus V\xrightarrow{(v,w)\mapsto v+w}X
 \xrightarrow{s\mapsto se_{10}-se_{00}}V.
\]
Quotient the acyclic first copy of $V$ in degree two
mapping identically onto $V\subset X$ in degree one.
The quotient is
$V[2]\oplus[X/V\to V]$, and the projection to the
trivial module uses only degree-zero augmentation.
After derived $N$ coinvariants and passage to the cone,
$H_3(Y)$ contains the direct summand $V_N\ne0$.

The only remaining core partition is two fixed points
and one regular two-point orbit, using all four states.
Here $S=U$ and the row actors agree, so the same
projective tensor complex has $d_1=0$.  Its degree-zero
term is $\mathbb Q[Se_{00}]$.  The degree-two map
\[
 (\mathbb Q[Se_{00}])^2\longrightarrow\mathbb Q[S]
\]
sends the first copy to the zero-column states and the
second to the corresponding one-column states.  It is
onto, and its kernel is the canonical copy of
$\mathbb Q[F]$ spanned by $(f,-f)$ for $f\in F$.
The kernel inclusion is $N$ equivariant, including when
an $N$ actor sends a regular orbit to a fixed point.
It is a quasi-isomorphism from $\mathbb Q[F][2]$ to
the positive-degree two-term complex.  Consequently the
whole complex is naturally quasi-isomorphic to
\[
 \mathbb Q[Se_{00}][0]\oplus\mathbb Q[F][2],
\]
with projection to $\mathbb Q$ on the first summand
only.  Again $H_3(Y)$ contains
$\mathbb Q[F]_N\ne0$.  Every case has been excluded.
\end{proof}

\begin{theorem}[Proved: the remaining five-center four-state star]
\label{CP2:A-star-five-seven-four}
Let $|M|=5$, $|N|\le7$ and $1\le|S|\le4$.  The
coned bridge $Y$ cannot have the homotopy type of
$\CP^2$.  In particular this excludes the original
eighteen-arrow star profile $10=(5,1,7,1,4)$.
\end{theorem}
\begin{proof}
The preceding lemma handles $M=P_5$.  Otherwise the
small-minimal-ideal comparison identifies $BM$ with the
classifying space of its minimal finite group.  If that
group is trivial, $BM$ is contractible.  Attaching the
additional singleton leaf cones off a contractible
corner, so $Y$ is homotopy equivalent to the original
two-object collage on $(M,N,S)$.  That collage has at
most $5+7+4=16$ arrows, contradicting the established
global lower bound.

If the minimal ideal has one row, the trivial left
$\mathbb QM$ module is projective.  The actual orbit
quotient gives a rational cohomology-ring equivalence
from $Y$ to the ordinary left $N$ act cone on $S/M$.
Corollary~\ref{CP2:A-high-seven-act} makes all its
degree-two products zero.  The only case left by the
five-element budget is
\[
 M=(\{0,1\}\times C_2)^1,\qquad
 (i,g)(j,h)=(i,gh).
\]
Indeed a nontrivial minimal group with at least two
rows already uses four elements, forcing two rows,
one column, group $C_2$, and only the global identity
outside the ideal.

Its core $U=SK$ is a disjoint union of one- or two-point
$C_2$ orbits.  If it has just one orbit, every commuting
$N$ actor permutes that orbit.  The covering
$B(N,U)\to BN$ factors through $E=B(N,S,M)$ using
empty $M$ strings.  Its rational cohomology pullback is
injective by transfer, so the restriction of every
second class of $Y$ to $BN$ is zero.  The coned-bridge
relative-cup lemma applies.  Likewise, if exactly one
core orbit is a singleton, it is the unique globally
right-fixed state, hence is $N$ fixed, and supplies a
section.

If $S=U$, every orbit module is projective as a right
$\mathbb QM$ module.  To see this, choose
$e=(0,1)$: the right ideal $eM=C_2$ has permutation
algebra $e\mathbb QM$, and every rational $C_2$ orbit
module is a summand of copies of this module by finite
group averaging.  Thus the actual orbit augmentation
$B(S,M)\to S/M$ is an $N$-equivariant rational homology
equivalence.  Derived $N$ coinvariants identify the
coned quotient with the ordinary $N$ act cone on $S/M$
in rational cohomology, through the actual orbit functor.
The unrestricted order-seven act theorem again kills
all second products.

The state bound leaves only the case in which every
core orbit is a singleton.  The right action then
factors through the monoid quotient
\[
 q:M\longrightarrow\overline M=\{0,1\}^{\mathrm{left},1},
 \qquad q(i,g)=i,
\]
because $s(i,g)=(s(i,1))(i,g)=s(i,1)$: the intermediate
core state is fixed.  Choosing $(i,1)$ gives a literal
monoid section, so the entire star on $\overline M$
is an actual retract of the original star, with exactly
two fewer arrows.

We check the required rational comparison on the actual
map.  In $R=\mathbb QM$ put
$p=((0,1)-(0,-1))/2$.  The kernel $J$ of
$R\to\mathbb Q\overline M$ is the principal left
ideal $Rp$, and $p^2=p$.  Thus
$0\to J\to R\to\mathbb Q\overline M\to0$ is a
left-projective resolution.  The factorized state module
has $\mathbb Q[S]p=0$, so tensoring kills $J$.
Change of rings gives the natural $N$-equivariant
comparison
\[
 \mathbb Q[S]\otimes_R^{\mathbf L}\mathbb Q
 \simeq
 \mathbb Q[S]\otimes_{\mathbb Q\overline M}^{\mathbf L}
 \mathbb Q.
\]
The quotient functor on the coned stars consequently
induces a rational homology equivalence, in particular
a second-cohomology isomorphism.  The actual retraction
then transfers projective-plane type integrally: split
integral homology forces $H_2=\mathbb Z$, naturality of
the primitive square forces $H_4=\mathbb Z$, and the
retracted fundamental group is trivial.  The smaller
star would be a projective-plane model with at most
sixteen arrows, again contradicting the global lower
bound.  This finishes the five-element classification.
\end{proof}
